\chapter{Technologies and Tools Used}
\label{ch:tecnologias}

\section{Introduction}

This chapter summarises the technologies used in \rat: backend, interfaces, sandbox, and report.

\section{Python and FastAPI}

\textbf{Python~3.10+} was adopted as the main backend language, due to its maturity in the
information security ecosystem and the wide availability of libraries for PE file analysis,
binary format parsing, and integration with decompilation tools.
\textbf{FastAPI} \cite{fastapi2024} was chosen as the web framework for providing automatic
data validation with Pydantic, interactive OpenAPI documentation (\texttt{/docs}), and native
support for asynchronous NDJSON streaming (\texttt{/api/analyze\_stream}), essential for
real-time feedback during analysis. There is also an optional \textbf{Tkinter} graphical interface
for local static analysis without a browser, useful for quick demonstrations.

\section{React, Vite and TypeScript}

The web interface was built with \textbf{React~18} and \textbf{TypeScript}, ensuring static
typing that reduces errors during development. The \textbf{Vite} bundler enables instant hot-reload
during development and optimised production builds. For visual styling,
\textbf{Tailwind CSS} and \textbf{shadcn/ui} were adopted; transition animations are handled
with \textbf{Framer Motion}. The unit test suite uses \textbf{Vitest} (62 unit tests).
Integration with \textbf{Google Gemini} (\texttt{gemini-2.0-flash}) is done directly in the
browser: the user's API key is stored in \texttt{localStorage} and never passes through the
FastAPI backend, ensuring the server has no access to third-party credentials.

\section{.NET~8, WPF and Minimal API}

Three .NET~8 projects share the \texttt{RatAnalyzer.sln} solution: \textbf{Rat Analyzer}
(WPF desktop, compiles to \texttt{Rat Analyzer.exe}), \textbf{VmAgent} (Minimal API running inside the sandbox
VM to receive and execute samples), and \textbf{BenignVmTest} (benign executable to validate
the dynamic pipeline without risk). WPF serves as an alternative entry point to the web
interface, with automatic dependency bootstrap, persistent language selection, and direct
integration with the Hyper-V sandbox (Path~B).

\section{YARA, Ghidra and ILSpy}

\textbf{YARA} \cite{yara2023} is the pattern-matching rule language used for signature-based
detection. Rules reside in \filepath{yara\_rules/} and are compiled and applied to the binary under
analysis. \textbf{ILSpy CLI} \cite{ilspy2024} decompiles .NET assemblies into readable C\#
code, enabling inspection of executables written with the managed runtime.
\textbf{Ghidra~12+} \cite{ghidra2019}, developed by the NSA, is used for decompilation of
native PE binaries, producing pseudo-C that is exposed in the web interface code panel.

\section{Hyper-V, PowerShell and SQLite}

The dynamic sandbox relies on \textbf{Hyper-V}, the native Windows hypervisor, with a VM named
\texttt{MalwareSandbox}, connected to an \textbf{Internal} switch with no Internet connectivity
on the guest and a \texttt{CleanState} snapshot that is restored before each analysis.
\textbf{PowerShell} automates one-time VM setup (\filepath{01-Setup-MalwareSandbox.ps1}) and,
on \textbf{Path~B}, sample copy, guest execution, and report transfer
(\filepath{04-Run-Sample.ps1}). On \textbf{Path~A}, the Python backend
(\filepath{vm\_drivers/hyperv.py}) uses PowerShell only for snapshot restore and VM start;
sample upload, execution, and report retrieval go through the HTTP vm-agent.
\textbf{SQLite} persists analysis jobs in \texttt{analysis.db}; \textbf{Redis}
with \textbf{RQ} is a configurable option for distributed queuing in environments with multiple workers.

The backend uses a pluggable sandbox driver so that VM orchestration can change between
development, Hyper-V, or future Proxmox deployments without rewriting the job pipeline. The
orchestrator calls the active driver for restore, execution, and report retrieval while keeping
the analysis contract unchanged.

\section{GitHub and LaTeX}

\textbf{GitHub} hosts the project repository, with Git version control and continuous integration
via \textbf{GitHub Actions} (CI with pytest, Vitest, and dotnet build). Changes are recorded in
\texttt{CHANGELOG.md} and development conventions in the root \texttt{README.md}. This report was
written in \textbf{\LaTeX}, a typesetting system widely used in academic and scientific contexts,
which automatically manages section numbering, figures, tables, and bibliographic references
with rigour and consistency.

\section{Security and Isolation}

The design adopts \emph{fail-closed} in production: the backend refuses to start without
\texttt{RATANALYZER\_API\_TOKEN} when \texttt{RATANALYZER\_ENV=production}; the vm-agent
refuses to start without \texttt{VM\_AGENT\_TOKEN}, except when
\texttt{VM\_AGENT\_ALLOW\_INSECURE=1} in local development.

Secrets reside in \filepath{secrets/} (gitignored) and are produced by
\filepath{scripts/generate-production-secrets.ps1}.
HTTP uploads require authentication, with name sanitisation, size limits, and storage in
\texttt{DATA\_DIR} outside the repository. The canonical security documentation is in
\filepath{docs/SEGURANCA.md}.

\section{Technology Choices and Rejected Alternatives}

Major stack decisions balanced \textbf{local reproducibility}, \textbf{Windows fidelity}, and
\textbf{academic maintainability}. Hyper-V with snapshot restore was preferred over Docker because
RAT behaviour depends on full Win32 semantics; Docker excels at service isolation but poorly
replicates GUI, registry, and injection patterns relevant to this project. KVM offers comparable
isolation on Linux hosts but would have forced a heterogeneous lab setup on predominantly Windows
workstations. FastAPI was chosen over Flask or Django for asynchronous streaming and schema
validation; React over alternative SPA frameworks for ecosystem familiarity and test tooling; WPF
over Electron to avoid bundling a browser runtime while retaining direct Hyper-V integration.
Redis-backed queuing remains optional so the platform operates offline with local worker threads.
The Gemini assistant runs in the browser only: decompiled excerpts and API keys stay on the client,
which protects sample confidentiality at the cost of centralised LLM control.

\section{Conclusion}

The stack reflects the heterogeneous nature of malware analysis: Python for pipeline and API,
TypeScript for UX, C\# for Windows/Hyper-V integration, and PowerShell for automation.
